\begin{abstract}
Medication reconciliation at clinical handoffs is a high-stakes, error-prone process.
Large language models are increasingly proposed to assist with this task using
FHIR-structured patient records, but a fundamental and largely unstudied variable is
how the FHIR data is serialised before being passed to the model. We present the first
systematic comparison of four FHIR serialisation strategies (Raw JSON, Markdown Table,
Clinical Narrative, and Chronological Timeline) across five open-weight models
(Phi-3.5-mini, Mistral-7B, BioMistral-7B, Llama-3.1-8B, Llama-3.3-70B) on a
controlled benchmark of 200 synthetic patients, totalling 4,000 inference runs. We
find that serialisation strategy has a large, statistically significant effect on
performance for models up to 8B parameters: Clinical Narrative outperforms Raw JSON
by up to 19 F1 points for Mistral-7B ($r = 0.617$, $p < 10^{-10}$). This advantage
reverses at 70B, where Raw JSON achieves the best mean F1 of 0.9956. In all 20
model and strategy combinations, mean precision exceeds mean recall: omission is the
dominant failure mode, with models more often missing an active medication than
fabricating one, which changes how clinical safety auditing priorities should be set. Smaller models plateau at roughly 7--10 concurrent
active medications, leaving polypharmacy patients, the patients most at risk from
reconciliation errors, systematically underserved. BioMistral-7B, a domain-pretrained
model without instruction tuning, produces zero usable output in all conditions,
showing that domain pretraining alone is not sufficient for structured extraction.
These results offer practical, evidence-based format recommendations for clinical LLM
deployment: Clinical Narrative for models up to 8B, Raw JSON for 70B and above. The
complete pipeline is reproducible on open-source tools running on an AWS
\texttt{g6e.xlarge} instance (NVIDIA L40S, 48 GB VRAM).
\end{abstract}
